%%%%%%%%%%%%%%%%%%%%%%%%%%%%%%%%%%%%%%%%%%%%%%%%%%%%%%%%%%%%%%%%%%%%%%%
% QE-CONVERSE user's manual
%%%%%%%%%%%%%%%%%%%%%%%%%%%%%%%%%%%%%%%%%%%%%%%%%%%%%%%%%%%%%%%%%%%%%%%
\documentclass[a4paper,11pt,twoside]{article}
\usepackage[utf8]{inputenc}             % For accents and weird symbols
%\usepackage{braket}                     % bra-ket's
\usepackage{amsmath}                    % extra mathematical symbols
\usepackage{bm}                         % bold mathematical symbols
\usepackage{verbatim}                   % insert verbatim files
\usepackage[hang,small,bf]{caption}     % customizable captions
\usepackage{psfrag}                     % smart postscript substitutions
\usepackage{fancyvrb}                   % for including source code
\usepackage[in]{fullpage}               % use all page
\usepackage[pdftex]{graphicx}           % to insert figures
\usepackage{epstopdf}                   % convert eps to pdf
\DeclareGraphicsExtensions{.pdf,.png,.jpg,.eps}
\usepackage{hyperref}                   % crossref in PDF
\hypersetup{pdftex, colorlinks=true, linkcolor=blue, citecolor=blue,
            filecolor=blue, urlcolor=blue, pdftitle=, pdfauthor=Davide Ceresoli,
            pdfsubject=QE-GIPAW user's manual, pdfkeywords=}

%===================================================================================================
\title{QE-CONVERSE user's manual}
\author{Simone Fioccola, \texttt{sfioccola@laas.fr}}
\date{Last revision: \today}
%===================================================================================================

%===================================================================================================
\begin{document}
%===================================================================================================
\maketitle
\thispagestyle{empty}
%\tableofcontents
%\clearpage

%===================================================================================================
\section{Introduction}
%===================================================================================================
The QE-CONVERSE implement a non-perturbative approach (converse) to compute the orbital magnetization in isolated and periodic systems. The calculation of orbital magnetization allows ab-initio computation of macroscopic properties like the Nuclear Magnetic Resonance (NMR) chemical shifts and the Electronic Paramagnetic Resonance (EPR) g tensor.

%===================================================================================================
\section{Features}
%===================================================================================================
\begin{itemize}
  \item Periodic and isolated systems
  \item Norm-conserving pseudopotentials
  \item Parallelization over bands and g-vectors
  \item NMR shielding tensors
  \item EPR g-tensor
  \item Electric field gradient (EFG) tensors and NMR/NQR quadrupolar parameters
  \item LDA and GGA functionals
\end{itemize}

%===================================================================================================
\section{Author contributions}
%===================================================================================================
{
  S. Fioccola, L. Giacomazzi, D. Ceresoli, N. Richard, A. Hemeryck, L. Martin-Samos
} 
%===================================================================================================
\section{Build instructions}
%===================================================================================================
This current version of the code is compatible with the version 7.5 of Quantum-Espresso package.
\begin{enumerate}
\item A Quantum-Espresso package version 7.5 must be previously installed (\url{https://gitlab.com/QEF/q-e/-/releases/qe-7.5}). To take advantage of the enhancements in linear algebra operations, the configuration with scaLAPACK package or ELPA library is suggested.
QE user's guide at \url{http://www.quantum-espresso.org/user_guide/user_guide.html}
\item Type:
\begin{verbatim}
    git clone https://github.com/mammasmias/QE-CONVERSE 
\end{verbatim}
This will download from github the latest stable version of QE-CONVERSE.
\item Type:
\begin{verbatim}
cd QE-CONVERSE
\end{verbatim}
\item 
 Copy the source files and the Makefile from \texttt{/src/} dictory into \texttt{/PP/} directory of Quantum-Espresso.
\item
 In the main directory of QE-7.5 type \texttt{make pp}. You should find the binary files \texttt{qe-converse.x} (NMR/EPR) and \texttt{qe-efg.x} (EFG) in the \texttt{/bin/} directory of Quantum-Espresso.
\end{enumerate}
%===================================================================================================
\section{Quick start}
%===================================================================================================
To calculate NMR/EPR parameters you need: 
\begin{enumerate}
\item pseudopotentials containing the GIPAW reconstruction (\url{https://sites.google.com/site/dceresoli/pseudopotentials})
\item run \texttt{pw.x} to perform the SCF calculation 
\item run \texttt{qe-converse.x} to calculate parameters (look into folder
\texttt{examples} for NMR shielding, EPR g-tensor.
)
\end{enumerate}

%===================================================================================================
\section{Input file description}
%===================================================================================================
The input file consists on only one namelist \texttt{\&input\textunderscore{}qeconverse} with the
following keywords:
\begin{description}

\item[prefix] (type: character, default: \texttt{'prefix'})\\
Description: prefix of files saved by program \texttt{pw.x}. The value of this keyword must be the same used in the SCF calculation.

\item[outdir] (type: character, default: \texttt{'./'})\\
Description: temporary directory for \texttt{pw.x} restart files. The value of this keyword must be the same used in the SCF calculation.

\item[diagonalization] (type: string, default: \texttt{'david'})\\
Description: diagonalization method (only allowed values: \texttt{'david'} )

\item[verbosity] (type: string default: \texttt{'high'})\\
Description: verbosity level (allowed values: \texttt{'low'}, \texttt{'medium'}, \texttt{'high'})

\item[q\_gipaw] (type: real, default: 0.01, units: bohrradius$^{-1}$)\\
Description: the small wave-vector for the covariant finite difference formula.

\item[dudk\_method] (type: string, default: \texttt{'covariant'})\\
Description: k-point derivative method: \texttt{'covariant'},
\texttt{'singlepoint'}, or \texttt{'kdotp'}.
The \texttt{'kdotp'} method solves the projected Sternheimer equation using
QE's conjugate-gradient solver and the full GIPAW Hamiltonian. Its Hamiltonian
derivative is evaluated by central differences with step \texttt{q\_gipaw};
step convergence must therefore still be checked. It uses the same occupied-band
subspace as \texttt{'covariant'}, including the occupation cutoff for smeared
calculations. It requires complex, collinear, norm-conserving wavefunctions.
The linear-system residual tolerance is the smaller of
$\sqrt{\mathtt{conv\_threshold}}$ and $10^{-8}$ (Rydberg units).
A failure of the linear solver terminates the calculation.

\item[diag\_thr\_init] (type: real, default: 10$^{-7}$, units: Ry$^2$)\\
Description: Convergence threshold (ethr) for iterative diagonalization.

\item[conv\_threshold] (type: real, default: 10$^{-8}$, units: Ry$^2$)\\
Description: convergence threshold for the diagonalization in the SCF step.

\item[mixing\_beta] (type: real, default: 0.5)\\
Description: mixing factor for self-consistency.

\item[lambda\_so(1,..,3)] (type: real, default: 0.0, units: Bohr magneton)\\
Description: Cartesian components of electron spin. The value (1,..,3) denotes the spin-orbit coupling direction.

\item[m\_0(1,..,3)] (type: real, default: 0.0, units: nuclear magneton)\\
Description: Cartesian components of nuclear dipole. The value (1,..,3) denotes the nuclear dipole moment direction.

\item[m\_0\_atom] (type: integer, default: 0)\\
Description: Atom index carrying the nuclear magnetic dipole.

\item[delete\_dudk\_files] (type: logical, default: \texttt{.false.})\\
Description: If \texttt{.true.}, temporary du/dk files (wavefunction k-derivatives) are deleted after the orbital magnetization calculation.

\item[dudk\_in\_memory] (type: logical, default: \texttt{.false.})\\
Description: If \texttt{.true.}, du/dk gradients (wavefunction k-derivatives) are stored in memory rather than written to disk. This is useful when there is enough memory to store these gradients.


\end{description}


%===================================================================================================
\section{Electric field gradient (EFG)}
%===================================================================================================
The electric field gradient at each nucleus is a ground-state property and is computed by a
separate binary, \texttt{qe-efg.x}, directly from the converged SCF charge density and
wavefunctions (no converse SCF is needed). A single ordinary \texttt{pw.x} run is sufficient;
symmetry may be left on. The total EFG is the sum of the bare valence, ionic (Ewald), and GIPAW
reconstruction contributions, which are printed separately together with the principal values and
the quadrupolar parameters $C_q$, $\eta$, and $\nu_Q$.

The input is the namelist \texttt{\&input\textunderscore{}qeefg}:
\begin{description}
\item[prefix] (character) prefix used in the SCF \texttt{pw.x} calculation.
\item[outdir] (character) scratch directory used in the SCF \texttt{pw.x} calculation.
\item[q\textunderscore{}efg($n$)] (real, default $0$) nuclear quadrupole moment of atom type $n$,
in units of $10^{-30}\,\mathrm{m}^2$ (= 10 mbarn); $C_q$ is printed only for types with a non-zero value.
\item[i\textunderscore{}efg($n$)] (real, default $0$) nuclear spin $I$ of atom type $n$; the
quadrupolar frequency $\nu_Q$ is printed only for types with $I \ge 1$.
\end{description}
Atom types are indexed as in the \texttt{ATOMIC\textunderscore{}SPECIES} block of the \texttt{pw.x}
input. Run \texttt{pw.x} followed by \texttt{qe-efg.x} to obtain the EFG at every site.

%===================================================================================================
\section{Limitations}
%===================================================================================================
K\_POINTS GAMMA are not supported. For Gamma-point calculations use: K\_POINTS automatic and `1 1 1 0 0 0`


%===================================================================================================
\section{Resources}
%===================================================================================================
\begin{itemize}

\item NMR periodic table: \url{http://www.pascal-man.com/periodic-table/periodictable.html}

\end{itemize}


%===================================================================================================
\begin{thebibliography}{99}
%===================================================================================================
\bibitem{}
  

\end{thebibliography}

%===================================================================================================
\end{document}
%===================================================================================================

